\pg{341} b) O \emph{infinito conhece-se effectivamente} pelo \emph{finito}. --- É certo que temos uma idea distincta do \emph{infinito}, a qual é \emph{positivo-negativa}, emquanto representa uma entidade não-limitada. Essa idea, quanto á parte \emph{negativa}, deriva --- ou \emph{do proprio infinito}, --- ou \emph{da nossa intelligencia}, --- ou \emph{do meio objectivo do conhecimento}. Ora não pode derivar --- nem \emph{do proprio infinito}, pois este é eminentemente positivo, --- nem \emph{da nossa intelligencia}, pois quanto mais positivo é o objecto cognoscivel, tanto mais positivo é o modo por que o conhecemos. Logo deve derivar do \emph{meio objectivo}, de que a intelligencia parte para chegar ao conhecimento do \emph{infinito} e que é o \emph{finito}, emquanto se exclue da sua entidade todo o limite. Logo o \emph{infinito} conhece-se effectivamente pelo \emph{finito} (').\footnote{Este modo de conhecer o \emph{infinito} é muito natural. O primeiro objecto do nosso conhecimento deve ser o que é mais proporcionado com a indole da nossa intelligencia. Ora a nossa intelligencia, sendo finita e limitada, deve primeiramente conhecer os objectos, que são finitos e limitados ; aliás não haveria proporção entre a mesma faculdade e o seu objecto connatural. Logo é do conhecimento do finito que devemos elevar-nos ao conhecimento do infinito. --- Alem d'isso, todo o nosso conhecimento intellectual é precedido pelo conhecimento sensitivo e deriva d'elle, emquanto nada chega á intelligencia, que não tenha passado pelos sentidos. Ora o objecto do conhecimento sensitivo é o ente visivel e finito. Logo o conhecimento do infinito deve ser precedido pelo conhecimento do finito, e deriva d' este.\par Portanto a nossa intelligencia forma a idea do infinito, emquanto exclue da entidade finita todos os limites, porque esta exclusão constitue a essencia do infinito e o extrema dos seres finitos. \emph{Spencer} escreve: «A idea que temos do limitado é composta da idea de uma determinada especie de ser, e da idea dos limites, em que o ser é conhecido. Na sua antithese, na concepção do illimitado, supprime-se a idea dos limites, mas permanece a idea do ser» (\emph{Premiers principes}, 1.ª partie, ch. IV). Muito bem dito. --- D' onde se vê o erro de \emph{Locke} e dos \emph{sensualistas}, quando, confundindo o \emph{infinito}\pgn{342}com o \emph{indefinito}, disseram que o \emph{infinito} se forma pela somma dos \emph{finitos}. O \emph{infinito} exclue todo o limite. Ora o finito sommado com o finito dá sempre o finito, e não o infinito. Alem d'isso, se o \emph{infinito} fosse o resultado da somma dos entes finitos, seguir-se-hia que, subtrahindo, uma a uma, as partes sommadas, deveriamos chegar a exgottar o infinito, que por isso deixaria de ser infinito, pois este é necessariamente inexgottavel.\par Os principaes erros, relativos á \emph{origem} da idea do \emph{infinito}, podem reduzir-se a tres : ao dos \emph{Cartesianos} e dos \emph{innatistas}, ao dos sequazes do \emph{ontologismo}, e ao dos \emph{pragmatistas}. --- Os \emph{Cartesianos} e os \emph{innatistas} sustentam que a idea do infinito é infundida por \textsc{Deus} na intelligencia humana, e que ella é o fundamento de todas as nossas percepções. Mas não é assim; porque, como veremos, não se dão ideas \emph{innatas}, e, para julgarmos que um ente é \emph{limitado}, não é necessario recorrer ao ente \emph{illimitado}, mas basta considerar as creaturas, cada uma das quaes, por não ter as perfeições das outras, concebe-se como limitada. --- Os \emph{ontologistas} ensinam que percebemos o \emph{infinito} em si mesmo, por intuição, e que da idea do infinito deduzimos a do finito. Tambem esta opinião é falsa; pois ninguem tem consciencia d' essa intuição, ou, melhor, todos temos consciencia de que não percebemos o \emph{infinito} por intuição, porque, se assim o percebessemos, a idea que temos d' elle seria perfeitissima, clarissima : o que não é. --- Os \emph{pragmatistas} sustentam com \emph{Hügel} que o homem tem uma experiencia intima, verdadeira e real, embora indeterminada, do \emph{infinito}, e que por isso pelo sentimento religioso e pela experiencia religiosa percebemos, ainda que de um modo confuso, o \emph{infinito}. Tambem esta opinião é inacceitavel. É certo que a nossa alma, creada por Deus e por Elle elevada á ordem sobrenatural, tem uma aspiração irresistivel para o infinito, um desejo ardente de conhecer perfeitamente a Deus, mas tambem é certo que o sentimento religioso e a experiencia religiosa, sendo coisas sujeitas a enganos, não são por si sufficientes para nos dar uma idéa exacta do infinito, e por isso precisam do auxilio do superior criterio da intelligencia, que as regule. Alem d'isso, só uma faculdade dotada de força abstractiva pode formar a idea do \emph{infinito}, excluindo da entidade finita todo o limite; ora o sentimento religioso e a consciencia religiosa não têm essa força abstractiva indispensavel. Finalmente, se o sentimento religioso e a consciencia religiosa podessem attingir ou perceber directamente o \emph{infinito}, a idea \emph{do infinito} seria exclusivamente \emph{positiva}; ora não é assim (\emph{Cf}. \emph{Enc}. « \emph{Pascendi} » do S. P. Pio X, 8 de Set. de 1907).}
